\documentclass[10pt,twocolumn]{article}
\usepackage[a4paper,margin=18mm,columnsep=6mm]{geometry}
\usepackage[T1]{fontenc}
\usepackage{lmodern}
\usepackage{graphicx}
\usepackage{booktabs}
\usepackage{caption}
\usepackage{microtype}
\usepackage[numbers,sort&compress]{natbib}
\usepackage[hidelinks]{hyperref}
\captionsetup{font=footnotesize,labelfont=bf}
\setlength{\bibsep}{1pt}
\setlength{\textfloatsep}{8pt plus 2pt minus 2pt}
\setlength{\floatsep}{6pt plus 2pt minus 2pt}
\newcommand{\rf}[1]{\texttt{\small #1}}

\title{What Survives a Screenshot? A quality matched robustness benchmark of invisible image watermarks against everyday sharing}
\author{Roie Zimlichman \and Noam Shaked \and Michael Varshay \and Gilad Kratter \and Rotem Cohen}
\date{Modern Computer Vision, final project, October 2026}

\begin{document}
\maketitle

\begin{abstract}
Prior benchmarks neglect the realistic transformations that occur during online sharing actions. We assess the classical DWT-DCT-SVD, HiDDeN, TrustMark, Watermark Anything, and PixelSeal architectures on 1,000 photographs and screen images under 75 sharing actions, including 15 real browser screenshots. To ensure a fair comparison independent of native output strengths or payload sizes, all embedders are calibrated to a uniform visual quality baseline of 40 dB PSNR. Our findings reveal that modern architectures (Watermark Anything and PixelSeal) survive all evaluated sharing workflows intact. In contrast, TrustMark fails only under crops, letterboxing and rotation, which realignment fully rescues, while the older methods degrade under lossy compression and resizing - HiDDeN mostly because the mark is destroyed, and DWT-DCT-SVD mostly because its block grid is misaligned.
\end{abstract}

\section{Introduction}
Invisible watermarks have become a standard mechanism for tracking the origin of AI-generated media. A watermark embeds a hidden message in an image, and a detector later tests a possibly altered copy for it. If a detector fails to recover the watermark, the failure stems from one of two distinct causes: either the watermark was completely destroyed within the pixel data, or the detector simply lost track of where the watermark sits due to a geometric misalignment. Out of the box, some models apply a heavy, easily detectable watermark (e.g., 33 dB PSNR), while others apply a faint, highly invisible one (e.g., 47 dB PSNR). A fair evaluation must first calibrate all models to the same visual volume.

\paragraph{Related work.}
While previous research establishes broad benchmarks focused on complex, deliberate removal attacks \cite{an2024wavesbenchmarkingrobustnessimage} and generative image editing \cite{lu2025wbench}, our work differs by exclusively targeting the unintentional transformations of everyday digital sharing. SyncSeal \cite{fernandez2025syncseal} shows that a learned synchroniser repairs geometric failures, but does not attribute failures per condition. We add screenshots captured from a real browser engine and explicitly diagnose whether watermark failures stem from data destruction or correctable geometric misalignment.
\section{Method}

\paragraph{Data.}
The evaluation dataset consists of 1,000 images, split evenly between 500 photographs and 500 screen captures.  The photographs are drawn from the professional and mobile sets of CLIC 2020 \cite{clic2020}. The screen images represent media commonly shared online, consisting of 250 application and web screenshots (from ScreenSpot-v2 \cite{wu2024osatlas}, ScreenSpot-Pro \cite{li2025screenspotpro}, and VisualWebBench \cite{liu2024visualwebbench}) and 250 document pages (from DocLayNet \cite{pfitzmann2022doclaynet} and arXiv). All images were resized to 1024 pixels on their longest side. To establish strict baseline detection thresholds, a separate null set of 5,000 unwatermarked images,  incorporating the same sources plus DIV2K \cite{agustsson2017div2k}, was additionally prepared. 

\paragraph{Watermarking methods.}
The benchmark evaluates five distinct post-hoc watermarking models: DWT-DCT-SVD \cite{invisiblewatermark}, TrustMark \cite{bui2023trustmark}, Watermark Anything \cite{sander2024wam}, PixelSeal \cite{fernandez2024videoseal, pixelseal2025}, and HiDDeN \cite{zhu2018hidden} via the 48-bit Stable Signature release \cite{fernandez2023stablesignature}.  Each model operates at its native payload, embedding a fresh random message into each image, 32, 100, 32, 256 and 48 bits in that order, at a native mean PSNR of 40.9, 42.8, 39.2, 46.7 and 32.6 dB (\rf{results/tables/quality\_full.tex}). TrustMark, Watermark Anything and PixelSeal decode a 256 px version of the image, while HiDDeN and DWT-DCT-SVD read the image at full size, the latter block by block.

\paragraph{Quality matching.}
Out of the box, the models alter images to differing degrees, producing native visibilities ranging from 33 dB PSNR (HiDDeN) to 47 dB PSNR (PixelSeal). To ensure a fair evaluation, every method was calibrated to a uniform visual volume of 40 dB PSNR per image prior to testing, the operating point also used by Evennou and Kijak \cite{evennou2026forensic}. For the four neural methods, the strength knob was adjusted directly until the image hit the 40 dB target. For the classical DWT-DCT-SVD method, its quantisation step was searched per image, and because finer steps stop decoding it reaches 39.0 dB on average. While matching is mathematically free for TrustMark, Watermark Anything, and PixelSeal, it causes HiDDeN to lose one bit in ten before any sharing actions are applied, clean bit accuracy 0.977 to 0.900 (\rf{results/tables/quality\_full.tex}).

\paragraph{Detection protocol.}
Performance is measured as the True Positive Rate (TPR) at a 1\% False Positive Rate (FPR), reported with 95\% Wilson intervals \cite{wilson1927}. The threshold is the 99th percentile of the null scores per method and per rung.

To distinguish whether a missing watermark was destroyed or merely misaligned, each image is decoded three times across distinct ``rungs''. Rung 1 decodes the shared image exactly as it is received. Rung 2 performs a perfect undo, filling regions lost to a crop with the unwatermarked cover. Rung 3 performs a practical undo, estimating the geometry from the received image and the original with DISK \cite{tyszkiewicz2020disk} and LightGlue \cite{lindenberger2023lightglue} keypoints, a RANSAC similarity fit and an ECC refinement~\cite{evangelidis2008ecc}, and uses its own threshold from realigned null images. Where a sharing action pushed more than 2 percent of null images over the threshold, the cell also carries a per family threshold. A mark lost at rung 1 but found at rung 2 is an alignment failure (the reader just lost track of where it sits), whereas a mark lost even at rung 2 has been completely destroyed.

\paragraph{Sharing actions.}
The pipeline tests 75 conditions grouped into 16 families, recording the exact parameters of each action and its inverse in order to reconstruct the image in rung 2. The set includes 45 single algorithmic transformations, such as JPEG compression, cropping, spatial shifts, letterboxing, and tone adjustments, and ten platform profiles that reproduce the resize and JPEG quality of Telegram, X, WhatsApp, Instagram and Facebook, on our images and on a 2048 px upscaled version of them. It adds 15 real browser screenshots generated via a headless Chromium engine \cite{playwright}. Finally, the pipeline evaluates five composite chains, which string multiple actions together to simulate multi-step sharing workflows.

\paragraph{Code and data.}
Code, data manifests, all results and all figures are in the public repository \cite{wmbench2026}. Paths below are relative to its root. Unless another file is named, every number in Section~\ref{sec:results} is a row of \rf{results/summary/cells\_full.csv}, which holds every condition and method at all three rungs with Wilson bounds.

\section{Results} \label{sec:results}

\begin{figure*}[t]
\centering
\includegraphics[width=0.66\textwidth]{figures/heatmap_rung1_two_blocks.png}
\caption{TPR at 1 percent FPR as received, 1,000 images at 40 dB, two of the four condition blocks, JPEG through shift above and the 15 real screenshots and the five chains below. A star marks a cell that also carries a per family threshold. The full map is \rf{results/summary/heatmap\_rung1\_full.png}, with rungs 2 and 3 next to it.}
\label{fig:heatmap}
\end{figure*}
\begin{figure*}[t]
\centering
\includegraphics[width=0.9\textwidth]{figures/curves_tpr_full_screenshot.png}
\caption{The 15 real browser screenshots. TPR at 1 percent FPR per condition, solid as received, dashed after the oracle realignment, shaded the 95 percent Wilson interval. Source: \rf{results/screenshot/curves\_tpr\_full.png}.}
\label{fig:screenshot}
\end{figure*}

Unless marked as smoke subset, all numbers in this section are from the full run on the 1,000 images. Table~\ref{tab:mean_TPR_table} gives the means and Figure~\ref{fig:heatmap} gives the cells.

DWT-DCT-SVD reads 0.00 to 0.02 under every downscale, upscale, crop, rotation of 3 degrees or more, shift of 8 pixels or more, and every real screenshot except the bit exact control, where it reads 0.85. The perfect undo lifts those cells to 0.60 to 0.85 and the practical undo to 0.47 to 0.83, the 30 percent crop excepted, which stays at 0.13 after the perfect undo.
HiDDeN reads 0.95 at crop 0.3 as received and 0.71 after the perfect undo, 0.48 and 0.31 on the phone to Instagram chain. 
No other method loses more than 0.04 in any cell after the undo. 
TrustMark averages 0.96, reads 1.00 under every JPEG, platform, screenshot, caption and sticker, and fails under crop 0.5 at 0.68, crop 0.3 at 0.09, the 15 percent letterbox at 0.12 and 10 degrees of rotation at 0.07. 
The oracle restores all four to 1.00 and the estimated realignment to 0.99 or 1.00. Watermark Anything and PixelSeal detect 1.00 of the images as received and never fall below 0.96 in a single cell. 


\paragraph{Screenshots.} 
Three methods score 1.00 under every one of the 15 real screenshots (Figure~\ref{fig:screenshot}). HiDDeN falls to 0.25 on the retina screen and returns to 0.97 after the undo. The classical method reads 0.01 under every screenshot and 0.78 to 0.83 after the undo, 0.60 when the screenshot is saved as JPEG. On the 50 image smoke subset, a browser free simulation of the same resampling and the real browser gave the same detection answer in 66 of the 75 screenshot, method and rung cells (\rf{results/calibration/screenshots\_smoke.csv}).

The practical undo registered 74,980 of 75,000 pairs with a median error of 0.13 px and trails the perfect undo by more than 0.1 in two
classical cells only, the quarter size downscale and the 20 percent caption band.

%הגרסה הקצרה יותר של headline_full
\begin{table}[t]
\centering\footnotesize
\setlength{\tabcolsep}{2.5pt}
\begin{tabular}{lrrrrr}
\toprule
Method & Rung 1 & Rung 2 & Rung 3 & Recov. & Destr. \\
\midrule
DWT-DCT-SVD & 0.26 & 0.68 & 0.66 & 0.41 & 0.32 \\
HiDDeN & 0.79 & 0.86 & 0.85 & 0.08 & 0.12 \\
TrustMark & 0.96 & 1.00 & 1.00 & 0.04 & 0.00 \\
Watermark Anything & 1.00 & 1.00 & 1.00 & 0.00 & 0.00 \\
PixelSeal & 1.00 & 1.00 & 1.00 & 0.00 & 0.00 \\
\bottomrule
\end{tabular}
\caption{Mean TPR at 1 percent FPR over the 75 conditions at the three rungs, and the mean share of images labelled recoverable, lost at rung 1 and found at rung 2, and destroyed, lost even at rung 2. The remainder, found at rung 2 but not at rung 3, is at most 0.02 for any method. Source: \rf{results/tables/headline\_full.tex}.}
\label{tab:mean_TPR_table}
\end{table}

\section{Discussion}
At equal visibility, the ranking is clear and stable across 75 conditions. Watermark Anything and PixelSeal survive every
everyday sharing action in this benchmark as received. TrustMark fails only under geometry that a realignment repairs
completely. HiDDeN loses mostly to destruction, under JPEG, the platform profiles and strong downscaling. DWT-DCT-SVD loses mostly to alignment, since any move of its block grid hides the mark, and the practical undo recovers most of it, 0.26 to 0.66 on average.
Realignment lowers HiDDeN under crops, the only case where a crop reads better on its own (\rf{results/crop/curves\_tpr\_full.png}). HiDDeN reads the whole image at full size, so after the perfect undo the unwatermarked cover pixels that fill the cropped region dilute the mark, while the other learned methods shrink the image before decoding.

\paragraph{Limitations.}
The screenshot conditions come from a real Chromium renderer validated against simulation, and screenshots from operating system tools, phones and messaging apps were not collected. The platform qualities for WhatsApp, Instagram and Facebook are estimates covered by a sensitivity sweep. All images are 1024 px on the long side, every method is compared at one quality point, and the classical method sits slightly below it. The rung 1 to rung 2 gap marks where a synchroniser such as SyncSeal \cite{fernandez2025syncseal} would pay, and evaluating one on exactly those cells is the natural next step.

\bibliographystyle{unsrt}
\scriptsize\bibliography{references}

\end{document}
